\documentclass{article}
\usepackage[utf8]{inputenc}
\usepackage[T1]{fontenc}
\usepackage{geometry}
\usepackage{graphicx}
\usepackage{float}
\usepackage{hyperref}
\usepackage[english,indonesian]{babel}

\title{Optimasi Pemrosesan Template LaTeX Menggunakan Pendekatan Berbasis AI}
\author{Academic Support Agent}
\date{\today}

\begin{document}
\maketitle

\begin{abstract}

\end{abstract}

\chapter{PENDAHULUAN}


\chapter{TINJAUAN PUSTAKA}


% Original content (commented out):
% \documentclass[a4paper,10pt]{article}
\usepackage{graphicx}
\usepackage{amsmath}
\usepackage[utf8]{inputenc}
\usepackage{hyperref}
\usepackage{indentfirst}
\setlength{\topmargin}{0pt}
\setlength{\headheight}{10pt}
\setlength{\headsep}{0.2cm}
\setlength{\oddsidemargin}{0.0cm}
\setlength{\textwidth}{160mm}
\setlength{\textheight}{235mm}
\setlength{\footskip}{15pt}
\setlength{\marginparwidth}{0pt}
\setlength{\marginparsep}{0pt}
\pagestyle{empty}


\headheight=0.9cm
%\usepackage{fancyhdr} %this packages provides fancy up and bottom of page
%\pagestyle{fancy}

%fancy
%\fancyhead[LO]{}
%\fancyhead[RO]{}
%\fancyfoot[C]{\footnotesize{3rd Russian School on Quantum %Technologies}}
%\fancyfoot[RO]{\thepage}
%\fancyfoot[CE]{\footnotesize{3rd Russian School on Quantum %Technologies}}
%\fancyfoot[LE]{ }

\usepackage{background}
\backgroundsetup{contents={\includegraphics[scale=1]{QTS23_logo.png}}, scale=0.3,placement=top,opacity=1, vshift=-40} 
%%%%%%%%%%%%%%%%%%%%%%%%%%%%%%%%%%%%%%%%%%%%%%%%%%%

\makeatother
\begin{document}
 
\begin{center}
\mbox{}

% TITLE
{\large{\bfseries  {Two-Page Summary Preparation for QTS. Important: Do Not Use Symbols, Special Characters, or Math in the title}}} % Please input your paper Title
\\[8pt]

% NAMES
\textbf{{\underline{Firstname1 Lastname1}$^{1*}$}, Firstname2 Lastname2$^2$, \\Firstname3 Lastname3$^{1,3}$}
\\[8pt]

% AFFILIATIONS
\small{\textit{
$^1$School of Energy and Power Engineering, Heat Transfer University, Xi'an, China \\
$^2$Institute of Power Engineering, Computational Science Research Center, Jeju, Korea \\
$^3$Department of Mechanical Engineering, University of Heat Transfer and Fluid Flow, Kyoto, Japan \\
}}

*E-mail: {\tt author@xx.yyy.zzz}

\end{center}
\vspace{0.2cm}

\begin{abstract}
This section of the paper should be a single paragraph abstract
outlining the aims, scope and conclusion of the summary. While no word limit is imposed, authors should aim for an abstract length of about 300 words for full papers or about 50 words for shorter contributions.    
\end{abstract}
\vspace{10pt}

Summaries should be submitted electronically via register form on {\tt \small http://qutes.org/} as original file in \LaTeX\ format, pdf file of your summary and files of figures, if they exists. You can submitted all files in archive (zip, rar or gz).


%%%%%%%%%%%%%%%%%%%%%%%%%%%%%%%%%%%%%%
%Figures
%%%%%%%%%%%%%%%%%%%%%%%%%%%%%%%%%%%%%%

 %	\begin{figure}[h]\label{fig1}
%		\centering\includegraphics[width=0.5\textwidth]{sample_figure.pdf}\caption{Caption}
%	\end{figure}

%%%%%%%%%%%%%%%%%%%%%%%%%%%%%%%%%%%%%%
%Equations
%%%%%%%%%%%%%%%%%%%%%%%%%%%%%%%%%%%%%%
\subsection*{Equations}

Equations should be placed on separate lines and numbered.
Examples of equations are given below.
\begin{equation}
q=-\lambda \left( \frac{\partial T}{\partial n} \right),
\label{eq1}
\end{equation}
where $\lambda$ is the thermal conductivity.
Equation~(\ref{eq1}) is the differential form of Fourier's Law of thermal conduction.

%%%%%%%%%%%%%%%%%%%%%%%%%%%%%%%%%%%%%%
%Tables
%%%%%%%%%%%%%%%%%%%%%%%%%%%%%%%%%%%%%%
\subsection*{Tables}

An example of a table is shown as Table~\ref{table1}.
Somewhat different styles are allowed according to the type and purpose of the table.
The caption text must be above the table.

\begin{table}[tbh]
\caption{\label{table1} This is an example of a table.}
\vspace{2mm}
\centerline{
\begin{tabular}{ccc}
\hline
format & size [mm$^2$] & aspect ratio\\
\hline
A4     & $210 \times 297$ & 1.414 \\
JIS B5 & $182 \times 257$ & 1.414 \\
Legal  & $216 \times 356$ & 1.647 \\
Letter & $216 \times 279$ & 1.294 \\
\hline
\end{tabular}}
\end{table}

%%%%%%%%%%%%%%%%%%%%%%%%%%%%%%%%%%%%%%
%Figures
%%%%%%%%%%%%%%%%%%%%%%%%%%%%%%%%%%%%%%

\subsection*{Figures}

Figures are to be centered, with the reference and caption printed
below the figure.

All figures must be quoted in the text. For example, see Fig.\,\ref{fig:1}.

If photographs or images are included, high quality originals should
be used. Figures should not appear any further from their first
citation in the text than is necessary for satisfactory layout.

\begin{figure}
  \centering
  \includegraphics[width=0.9\textwidth]{figures/figure-template.pdf}
  \caption{Entire caption for figure centered and below the illustration\label{fig:1}}
\end{figure}

%%%%%%%%%%%%%%%%%%%%%%%%%%%%%%%%%%%%%%
%References
%%%%%%%%%%%%%%%%%%%%%%%%%%%%%%%%%%%%%%
\subsection*{References}

References should be numbered in order of appearance, for example \cite{bib1}, \cite{bib2}, and \cite{bib3,bib4}.

%%%%%%%%%%%%%%%%%%%%%%%%%%%%%%%%%%%%%%
%acknowledgement
%%%%%%%%%%%%%%%%%%%%%%%%%%%%%%%%%%%%%%
\vspace{10pt}

Acknowledgements belong here.



\vspace{0.5cm}

\small{
\begin{thebibliography}{}
\bibitem{bib1}
\textit{O.L. Muskens, V. Giannini, J.A. S\'anchez-Gil, and J.G. Rivas},
Strong enhancement of the radiative decay rate of emitters by single plasmonic nanoantennas.
Nano Lett. {\bf7}, 2871 (2007).

\bibitem{bib2}
\textit{J.N. Farahani, D.W. Pohl, H.J. Eisler, and B. Hecht},
Single quantum dot coupled to a scanning optical antenna: A tunable superemitter.
Phys. Rev. Lett. {\bf95}, 017402 (2005).

\bibitem{bib3}
\textit{P.J. Schuck, D.P. Fromm, A. Sundaramurthy, G.S. Kino, and W.E. Moerner},
Improving the mismatch between light and nanoscale objects with gold bowtie nanoantennas.
Phys. Rev. Lett. {\bf94}, 017402 (2005).

\bibitem{bib4}
\textit{T.H. Taminiau, R.J. Moerland, F.B. Segerink, L. Kuipers, and N.F. van Hulst},
$\lambda/4$ resonance of an optical monopole antenna probed by single molecule fluorescence.
Nano Lett. {\bf7}, 28 (2007).

\end{thebibliography}}

\end{document}




\end{document}